\section{The Modern Magisterial File}

Four documents and one synod stand between the Second Vatican Council and the present law. Each belongs to a different genre and carries a different weight, and the fastest way to misuse the file is to quote them all as though they were the same kind of act.

\subsection{The Council: what it said and what it left}

\emph{Presbyterorum ordinis} 16 has already supplied this article's controlling conciliar clause---celibacy \latin{non exigitur quidem a sacerdotio suapte natura}---and the sentence in which the Council both refused to change the Eastern discipline and confirmed the Latin one. Two further features of that article are worth marking now that the codes have been read.

The Council states the Eastern position with unusual precision: in the Eastern traditions, ``besides those who with all the bishops, by a gift of grace, choose to observe celibacy, there are also married priests of highest merit,'' \latin{sunt etiam optime meriti Presbyteri coniugati}. The phrase \latin{cum omnibus Episcopis} concedes the unmarried episcopate as a fact of Eastern practice, which is what CCEO canon 180 3° would later state as a requirement of the bond. And the Council's account of how the Latin rule arose is a historical claim in a doctrinal document: celibacy ``first was recommended to priests, later in the Latin Church was imposed upon all who were to be promoted to sacred orders''---\latin{postea in Ecclesia Latina omnibus ad Ordinem sacrum promovendis lege impositus est}. Section~10 tests that claim against the acts.

\emph{Optatam totius} 10, on priestly formation, contains a phrase easy to read past: it addresses students ``who follow the venerable tradition of celibacy according to the holy and firm laws of their own rite.'' And its central sentence about how the state is to be received turns on a distinction this article has kept throughout: the students are to feel how gratefully the state is to be undertaken \latin{non quidem solum ut lege ecclesiastica praeceptus, sed ut pretiosum donum Dei humiliter impetrandum}---not only as commanded by ecclesiastical law, but as a precious gift of God to be humbly begged for.\endnote{\emph{Optatam totius} 10 (28 October 1965), Latin and English from the Holy See's web deliveries, fetched, hashed, and read 2026-07-25; \emph{AAS} 58 (1966) 713--727, n.~10 at p.~719. The Latin's opening restriction, \latin{secundum proprii ritus sanctas firmasque leges}, is rendered in the English delivery as ``according to the holy and fixed laws of their own rite.'' The Decree's own formula---ecclesiastical law \emph{and} gift---is the exact pairing that CIC c.~277 §1 would carry into the Code with \latin{peculiare Dei donum}.} That the same sentence calls it both a command of ecclesiastical law and a gift to be prayed for is not a rhetorical hedge; it is the Council's whole position on the matter in eleven words.

What the Council did not do was debate the discipline on the floor. Paul~VI reserved the question to himself by a letter to Cardinal Tisserant, Dean of the College of Cardinals, dated 10 October 1965 and read in the general congregation the following day. That letter is officially attested---the encyclical of 1967 cites it in its own first footnote---but its text is not published in the \emph{Acta Apostolicae Sedis} and could not be obtained at any Holy See delivery for this article, so nothing is quoted from it here and no characterization of its terms is offered.\endnote{Publication-local negative result, bounded. \emph{Sacerdotalis caelibatus} n.~2, note~1, in both the Holy See's English and Latin deliveries and in the \emph{Acta} print at \emph{AAS} 59 (1967) 659: \latin{Cf.\ Epistula die x mensis Oct.\ anno MDCCCCLXV ad Eminentissimum virum Eugenium Tisserant Card.\ data, posteroque die in generali Concilii congregatione lecta}---the Vatican web Latin delivery prints \emph{eirum} for \emph{virum}, a delivery defect recorded here. The Holy See's index of Paul~VI's 1965 letters lists one letter to Cardinal Tisserant, \emph{Laeto animo} of 9 November 1965, on the date of the Council's closing, and a search of the whole of \emph{AAS} 57 (1965) located no other. The October letter's text belongs to the \emph{Acta Synodalia} of the Council, which the Holy See has not digitized. What the encyclical does say about the reservation is only that Paul~VI had promised the Council Fathers ``to give new luster and strength to priestly celibacy in the world of today.'' The frequently repeated date of 11 October is the date the letter was read, not its date.}

\subsection{\emph{Sacerdotalis caelibatus} (1967)}

Paul~VI's encyclical of 24 June 1967 is the longest single magisterial treatment of the subject and the one most often quoted in fragments. Its structure is worth having whole, because the parts are load-bearing on each other.

It opens by naming the objections, which is unusual for the genre, and then declines to rest the law on necessity: n.~17 quotes the Council's \latin{non exigitur} against itself and says that the Council nonetheless ``did not hesitate to confirm solemnly the ancient, sacred and providential present law of priestly celibacy.'' The three sections that follow set out the reasons of fittingness as christological (nn.~19--25), ecclesiological (nn.~26--32), and eschatological (nn.~33--34): configuration to Christ, who ``remained throughout His whole life in the state of celibacy'' (n.~21); an undivided pastoral charity, celibacy being held ``as a symbol of, and stimulus to, charity'' (n.~24); and the sign of the world to come, in which ``they neither marry nor are given in marriage'' (n.~34). None is offered as a demonstration that a married priesthood is impossible; they are \latin{rationes convenientiae}, which is what \latin{non exigitur} leaves room for.

Two of its passages on the East deserve exact reading, because they are quoted in both directions. Number 38 attributes the Eastern discipline to a different historical background ``which the Holy Spirit has providentially and supernaturally influenced,'' and expresses esteem for the Eastern clergy. Number 40 observes that ``in the East only celibate priests are ordained bishops, and priests themselves cannot contract marriage after their ordination,'' concluding that those Churches ``also possess to a certain extent the principle of a celibate priesthood.'' The observation is accurate as to the two facts, and the inference is the encyclical's own; a reader may weigh it against CCEO canon 373's parallel honouring of both states.

And n.~42, treated in section~11, is the provision that made every subsequent married-convert-priest arrangement possible. It is worth noticing where it stands: in an encyclical whose purpose was to reaffirm the law, immediately after the section defending the Western tradition. The same document that says the law ``remains firm in its own right'' provides the instrument for excepting from it.

\subsection{\emph{Pastores dabo vobis} (1992)}

John Paul II's post-synodal exhortation on priestly formation returns to the terms this article began with, and states their relation directly: ``In virginity and celibacy, chastity retains its original meaning, that is, of human sexuality lived as a genuine sign of and precious service to the love of communion and gift of self to others.'' Celibacy is presented there as a mode of chastity, not as a substitute for it or an intensification of a different virtue.

Number 29 also records the synod fathers' own proposal in full, and the proposal is careful in a way worth noting: ``While in no way interfering with the discipline of the Oriental churches, the synod, in the conviction that perfect chastity in priestly celibacy is a charism, reminds priests that celibacy is a priceless gift of God for the Church\ldots\ This synod strongly reaffirms what the Latin Church and some Oriental rites require, that is, that the priesthood be conferred only on those men who have received from God the gift of the vocation to celibate chastity (without prejudice to the tradition of some Oriental churches and particular cases of married clergy who convert to Catholicism, which are admitted as exceptions in Pope Paul VI's encyclical on priestly celibacy, no.~42).''\endnote{\emph{Pastores dabo vobis} (25 March 1992) nn.~29 and 50, quoted from the Holy See's English web delivery, fetched, hashed, and read 2026-07-25; the Latin (\emph{AAS} 84 [1992]) was not collated. Two features of the quoted proposition are load-bearing for this article: it says ``the Latin Church \emph{and some Oriental rites},'' conceding that celibate discipline is not exclusively Latin; and it identifies the convert-clergy exception by its exact instrument, \emph{Sacerdotalis caelibatus} n.~42. Number 50 adds the formation principle that celibacy ``should not be considered just as a legal norm or as a totally external condition for admission to ordination.''} Both the ``some Oriental rites'' and the citation of \emph{Sacerdotalis caelibatus} 42 show a synod that knew exactly which distinctions it was making.

\subsection{The Amazon synod (2019) and \emph{Querida Amazonia} (2020)}

The Special Assembly of the Synod of Bishops for the Pan-Amazon Region, meeting 6--27 October 2019, addressed a shortage so severe that communities went months or years without the Eucharist. Its \emph{Final Document} n.~111 begins by affirming celibacy---``we value celibacy as a gift of God\ldots\ and we pray that there may be many vocations living the celibate priesthood''---and quotes \emph{Presbyterorum ordinis} 16's \latin{non exigitur} in Spanish translation before making its proposal:

\begin{studyframe}
\emph{proponemos establecer criterios y disposiciones de parte de la autoridad competente, en el marco de la Lumen Gentium 26, de ordenar sacerdotes a hombres idóneos y reconocidos de la comunidad, que tengan un diaconado permanente fecundo y reciban una formación adecuada para el presbiterado, pudiendo tener familia legítimamente constituída y estable, para sostener la vida de la comunidad cristiana mediante la predicación de la Palabra y la celebración de los Sacramentos en las zonas más remotas de la región amazónica.}

\emph{Working gloss:} we propose that the competent authority establish criteria and dispositions, within the framework of \emph{Lumen gentium} 26, for ordaining as priests suitable and recognized men of the community who have a fruitful permanent diaconate and receive an adequate formation for the presbyterate, who may have a legitimately constituted and stable family, in order to sustain the life of the Christian community through the preaching of the Word and the celebration of the Sacraments in the most remote areas of the Amazon region.
\end{studyframe}

The document adds, in one sentence, that ``in this regard, some declared themselves for a universal approach to the subject.''\endnote{Synod of Bishops, Special Assembly for the Pan-Amazon Region, \emph{Documento final}, n.~111, quoted from the Holy See Press Office daily bulletin B0820 of 26 October 2019, which publishes the document's Spanish text; fetched, hashed, and read 2026-07-25. There is no English text in that bulletin, and the document was not published in \emph{Acta Apostolicae Sedis}. The Spanish prints \emph{constituída} with that accent. The phrase \latin{viri probati}, universally used in reporting, does not occur in n.~111. On genre: a synodal final document is a consultative act of an assembly presented to the Roman Pontiff; it is not itself an act of the papal magisterium and enacts nothing. Number 110, immediately preceding, grounds the request in ``a right of the community to the celebration'' deriving from the essence of the Eucharist.} Two things about the proposal are frequently misreported. It asks for ordination of men who already have a fruitful permanent diaconate, not for the ordination of married men in general; and it asks the ``competent authority'' to establish criteria, which is a request for an act of governance, not a claim to perform one.

Pope Francis's post-synodal apostolic exhortation \emph{Querida Amazonia} of 2 February 2020 does not adopt the proposal, does not reject it, and does not mention it. This is a verifiable fact about the text rather than an impression: the official English text contains no occurrence of ``celibacy,'' ``celibate,'' ``viri probati,'' ``married men,'' ``ordain,'' or ``ordination''; the only cognate anywhere in it is the phrase ``ordained ministers,'' once, at n.~93.\endnote{\emph{Querida Amazonia} (2 February 2020), official English delivery, fetched, hashed, and read 2026-07-25; the searches reported are case-insensitive literal searches over the text extracted from that exact response. Bounded negative result: it establishes absence of those strings in that delivery of that translation, which for terms as basic as ``celibacy'' and ``ordination'' is strong evidence of the text's silence, but it is a search over one language of one delivery and is correctable. The exhortation's own account of its relation to the \emph{Final Document} is at nn.~2--3: ``I will not go into all of the issues treated at length in the final document. Nor do I claim to replace that text or to duplicate it''; ``I would like to officially present the Final Document\ldots\ I have preferred not to cite the Final Document in this Exhortation, because I would encourage everyone to read it in full.''}

What the exhortation does instead is answer the underlying problem by other routes: it insists that ``the exclusive character received in Holy Orders qualifies the priest alone to preside at the Eucharist,'' ``his particular, principal and non-delegable function'' (n.~87); it says that ``in the specific circumstances of the Amazon region\ldots\ a way must be found to ensure this priestly ministry'' (n.~89); it urges bishops ``especially those in Latin America'' to be ``more generous in encouraging those who display a missionary vocation to opt for the Amazon region'' (n.~90); it calls for many more permanent deacons and expanded lay responsibility (n.~92); and it warns that ``it is not simply a question of facilitating a greater presence of ordained ministers who can celebrate the Eucharist. That would be a very narrow aim'' (n.~93).

A reader may draw prudential conclusions from that silence, and many have drawn opposite ones. What cannot be drawn from it is a magisterial act: the exhortation neither changed the law nor pronounced on the proposal, and the law of section~3 through section~6 stands exactly as it stood in 2019.
